
\section{Method}

\subsection{Overview}

GQA is a \textbf{structural modification} of the attention mechanism: it reduces 
the number of $K, V$ heads from $H$ to $G$ while keeping all $H$ query heads 
intact. The modification consists of four components:

\begin{enumerate}
\item \textbf{Group configuration} ( 4.2): choose $G$.
\item \textbf{Projection initialization} ( 4.3): initialize $K, V$ for the reduced head count.
\item \textbf{Weight transfer} ( 4.4): copy query weights verbatim; group-reduce KV weights.
\item \textbf{Forward pass} ( 4.5): apply standard attention with KV repetition.

\end{enumerate}
\textbf{Figure 4.1} shows the difference between MHA, GQA, and MQA at $H = 8$.

\begin{lstlisting}
    MHA (H_KV = 8)          GQA (H_KV = 4)          MQA (H_KV = 1)

    Q1 -- K1 V1             Q1 -- K1 V1             Q1 -- K1 V1
    Q2 -- K2 V2             Q2 -- +                 Q2 -- +
    Q3 -- K3 V3             Q3 -- K2 V2             Q3 -- +
    Q4 -- K4 V4             Q4 -- +                 Q4 -- +
    Q5 -- K5 V5             Q5 -- K3 V3             Q5 -- +
    Q6 -- K6 V6             Q6 -- +                 Q6 -- +
    Q7 -- K7 V7             Q7 -- K4 V4             Q7 -- +
    Q8 -- K8 V8             Q8 -- +                 Q8 -- +
\end{lstlisting}

\subsection{Group Configuration}

The number of KV heads $G$ satisfies $1 \le G \le H$. Two natural choices:

\begin{itemize}
\item \textbf{Default:} $G = H / 2$ (2$$\times$ KV cache reduction).
\item \textbf{Aggressive:} $G = H / 4$ or $G = 1$ (MQA).

\end{itemize}
Our default is $G = H / 2$ to balance quality and memory.

\subsection{Projection Initialization}

For a target model with hidden size $d$ and per-head dimension $d_h$:

\[
W_Q \in \mathbb{R}^{d \times H d_h}, \qquad
W_K, W_V \in \mathbb{R}^{d \times G d_h}, \qquad
W_O \in \mathbb{R}^{H d_h \times d}
\tag{4.1}
\]

The query and output projections have unchanged shapes; only $W_K, W_V$ 
shrink by $H / G$.

\subsection{Weight Transfer}

\subsubsection{From MHA to GQA}

Given MHA weights $W_Q^{(\text{MHA})} \in \mathbb{R}^{d \times H d_h}$ and 
$W_K^{(\text{MHA})}, W_V^{(\text{MHA})} \in \mathbb{R}^{d \times H d_h}$, we 
initialize GQA weights as:

\[
W_Q^{(\text{GQA})} = W_Q^{(\text{MHA})}
\tag{4.2}
\]

\[
W_K^{(\text{GQA})}_g = \text{mean}\!\left(W_K^{(\text{MHA})}_{h \in g}\right), \quad
W_V^{(\text{GQA})}_g = \text{mean}\!\left(W_V^{(\text{MHA})}_{h \in g}\right)
\tag{4.3}
\]

where the mean is over the heads in group $g$. This is a \textbf{uniform average} 
of the head weights within each group.

\textbf{Alternative:} We could pick the first head in each group or use learned 
pooling. Our experiments use averaging.

\subsubsection{From GQA to GQA (Re-grouping)}

If the source is already GQA with $G_\text{src}$ groups and we want 
$G_\text{tgt} < G_\text{src}$, we group the source KV heads:

\[
W_K^{(\text{tgt})}_g = \text{mean}\!\left(W_K^{(\text{src})}_{g' \in g}\right)
\tag{4.4}
\]

If $G_\text{tgt} = G_\text{src}$, the copy is exact.

\subsubsection{Copy Procedure Pseudocode}

\textbf{Algorithm 1:} \texttt{copy\_weights} for GQA
\begin{lstlisting}
Input:  src_attn (MHA or GQA)
        H, G, d_h
Output: GQA weights

1:  # Query weights are always copied verbatim
2:  W_Q $\leftarrow$ extract_q(src_attn)
3:  copy(W_Q_gqa, W_Q)
4:
5:  # KV weights are grouped
6:  W_K_src $\leftarrow$ extract_k(src_attn)   # shape (d, H * d_h) for MHA
7:  W_V_src $\leftarrow$ extract_v(src_attn)
8:
9:  H_src $\leftarrow$ W_K_src.shape[1] // d_h
10: r $\leftarrow$ H_src // G
11:
12: for g in 0..G-1 do
13:     heads_in_g $\leftarrow$ src heads [g*r, (g+1)*r)
14:     W_K_gqa[g] $\leftarrow$ mean(W_K_src[heads_in_g])
15:     W_V_gqa[g] $\leftarrow$ mean(W_V_src[heads_in_g])
16: end for
17:
18: # Output projection unchanged
19: W_O_gqa $\leftarrow$ extract_o(src_attn)
20:
21: copy all projections into modules
22: freeze_base_weights()
\end{lstlisting}

\subsection{Forward Pass}

\textbf{Algorithm 2:} GQA forward
\begin{lstlisting}
Input:  X $\in$ R^{B $$\times$ N $$\times$ d}
        attention_mask
Output: Y $\in$ R^{B $$\times$ N $$\times$ d}

1:  Q $\leftarrow$ X W_Q            # (B, N, H, d_h)
2:  K $\leftarrow$ X W_K            # (B, N, G, d_h)
3:  V $\leftarrow$ X W_V            # (B, N, G, d_h)
4:
5:  if RoPE: (Q, K) $\leftarrow$ apply_rotary(Q, K)
6:
7:  # Repeat KV heads to match Q heads
8:  r $\leftarrow$ H // G
9:  K $\leftarrow$ repeat_kv(K, r)     # (B, N, H, d_h)
10: V $\leftarrow$ repeat_kv(V, r)
11:
12: # Standard scaled dot-product attention
13: scores $\leftarrow$ Q K^T / $\sqrt{}$d_h
14: scores $\leftarrow$ apply_causal_mask(scores)   # if decoder
15: scores $\leftarrow$ apply_padding_mask(scores)
16: attn $\leftarrow$ softmax(scores)
17: context $\leftarrow$ attn V
18:
19: Y $\leftarrow$ out_proj(context)
20: return Y
\end{lstlisting}

\subsection{KV Cache Analysis}

\subsubsection{Per-Token Memory}

For a single layer at sequence position $t$:

\[
\text{KV}^{(t)}_\text{MHA} = 2 H d_h \cdot \text{bytes}
\tag{4.5}
\]

\[
\text{KV}^{(t)}_\text{GQA} = 2 G d_h \cdot \text{bytes} = \frac{H}{G} \cdot \text{KV}^{(t)}_\text{MHA}
\tag{4.6}
\]

\subsubsection{Total Cache}

For sequence length $L$ and batch size $B$ across $N_l$ layers:

\[
\text{KV}_\text{total} = N_l \cdot L \cdot B \cdot 2 G d_h \cdot \text{bytes}
\tag{4.7}
\]

\subsubsection{Concrete Numbers}

For GPT-2 ($N_l = 12$, $H = 12$, $d_h = 64$, FP16):

\begin{table}[H]
\centering
\small
\begin{tabular}{lllll}
\toprule
\textbf{Config} & \textbf{$H_{KV}$} & \textbf{KV/token (KB)} & \textbf{KV at L=1024 (MB)} & \textbf{KV at L=8192 (MB)} \\
\midrule
MHA & 12 & 36.9 & 37.7 & 302 \\
GQA (H/2) & 6 & 18.4 & 18.9 & 151 \\
GQA (H/4) & 3 & 9.2 & 9.4 & 75 \\
MQA & 1 & 3.1 & 3.1 & 25 \\
\bottomrule
\end{tabular}
\end{table}

\subsection{Complexity Analysis}

\begin{table}[H]
\centering
\small
\begin{tabular}{lll}
\toprule
\textbf{Component} & \textbf{MHA} & \textbf{GQA} \\
\midrule
Q projection FLOPs & $\mathcal{O}(N d \cdot H d_h)$ & Same \\
K projection FLOPs & $\mathcal{O}(N d \cdot H d_h)$ & $\mathcal{O}(N d \cdot G d_h)$ \\
V projection FLOPs & Same as K & Same as K \\
Attention FLOPs & $\mathcal{O}(N^2 H d_h)$ & Same \\
Output projection FLOPs & Same & Same \\
\textbf{KV cache bytes} & $2 H d_h L B$ & $2 G d_h L B$ \\
\bottomrule
\end{tabular}
\end{table}

\textbf{Key insight.} GQA reduces KV cache and $K, V$ projection FLOPs proportionally 
to $G / H$, but \textbf{not} attention FLOPs or output projection FLOPs. The primary 
benefit is memory-bound inference.

\subsection{Rank and Head Grouping}

The number of KV heads $G$ must satisfy two constraints:

\begin{enumerate}
\item $G \le H$ (can't have more KV heads than query heads).
\item $H$ is divisible by $G$ (each group has $H / G$ heads).

\end{enumerate}
If $H$ is prime, only $G = 1$ or $G = H$ are valid. For $H = 12$, valid values 
are $G \in \{1, 2, 3, 4, 6, 12\}$.

\subsection{Theoretical Notes}

\subsubsection{MHA and MQA as Special Cases}

At $G = H$, GQA reduces to MHA (each head has its own KV). At $G = 1$, GQA 
reduces to MQA (all heads share KV). GQA is a \textbf{continuous interpolation} 
between these extremes.

\subsubsection{Expressiveness Bound}

For a fixed sequence, the set of attention matrices expressible by GQA is a 
subset of MHA's expressible set. The reduction is bounded by the number of 
distinct KV heads:

\[
|\text{Expr}(\text{GQA})| \le \binom{H}{G}
\tag{4.8}
\]

relative to MHA's $|\text{Expr}(\text{MHA})|$ (informally, up to a 
factor depending on $d_h$).